Deep Neural Networks achieve state-of-the-art predictive performance in cancer pharmacogenomics,
particularly for drug sensitivity prediction using high-dimensional genomic data. However, their
opaque black-box nature limits clinical adoption: clinicians cannot easily verify whether predictions
are grounded in genuine oncogenic biology or merely exploit learned statistical shortcuts, a
vulnerability known as the Clever Hans effect. This study proposes a Red Explainable Artificial
Intelligence (XAI) diagnostic pipeline that treats post-hoc explainability as an active model
interrogation tool rather than a justification mechanism. Using the Genomics of Drug Sensitivity in
Cancer (GDSC) dataset and the biologically informed DeepCCDS architecture, the pipeline applies
Sharpness-Aware Minimisation (SAM) to steer the model toward flatter loss landscape minima. Symbolic
if-then rules are extracted using a single-objective Genetic Algorithm and a multi-objective
Non-Dominated Sorting Genetic Algorithm II (NSGA-II), and evaluated for fidelity, comprehensibility,
and stability through a unified PROMETHEE II analysis, non-parametric Mann--Whitney U testing, and Gene
Set Enrichment Analysis (GSEA). The study hypothesises that landscape-aware regularisation will reduce
the generalisation gap and yield rules that are statistically simpler, more faithful to the model's
decision logic, and more strongly enriched in known oncogenic pathways than those from an unregularised
baseline.